\subsection*{More Questions}

Like C++ functions, Croqtile functions can be defined throughout a source file. When linking a module that contains one, its \textbf{function signature} determines how it is called.

A Croqtile function follows the C++ calling convention and uses C++ name mangling. It does not support \texttt{extern "C"}, because its spanned-data parameters use C++ template-based types.

Some features are not supported yet:

\begin{itemize}[leftmargin=1.6em]
\item Croqtile function declaration without definition. This allows calling Croqtile function from another module.
\item Croqtile function with \texttt{inline}. This allows the Croqtile function to be used in headers.
\end{itemize}
